% This content is included within ReaderDetailedPage; color is inherited.
% Use the available full width or minipage columns. Keep one page readable;
% additional pages are allowed by the helper, but must be reported accurately.
\textbf{1. 符号。}\(X\)是随机变量，\(p\)是正面出现的概率，\(\Pr(X=1)\)表示事件“\(X\)取值为一”的概率。这里\(\Pr\)不是额外的随机变量。
\par\medskip
\textbf{2. 公式含义。}期望是按概率加权的平均值：
\[
  \mathbb{E}[X]=\sum_{x\in\{0,1\}}x\Pr(X=x)
                =0(1-p)+1p=p.
\]
\textbf{3. 数值例子。}若\(p=0.7\)，期望为\(0.7\)。一次抛掷的取值仍只有零或一；期望描述大量重复实验的平均行为。
\par\medskip
\textbf{4. 来源边界。}这是模板自带的原创教学示例。实际精读中应明确标注哪些内容来自论文，哪些属于辅助推导。若怀疑原式有排印问题，应保留正文原式并在解释中给出理由。
\par\medskip
\textbf{5. 插页与页码。}本页页眉为“详细解释页”，所有内容为深红色，不显示页码。下一张正文页继续从正文上一页的下一页编号，物理页与正文页码的对应关系写入编译日志。
